\section*{4. 一般情形}

我们概述另一种证明 2.5 的方法, 它在任意特征下都成立, 但需假设以下陈述; 此陈述与 3.1 相近, 并是 XV, 1.1.4 的特例.

\medskip
\noindent(4.1) \underline{$\mathbb P(N)$ 中由满足 $x$ 是 $X\cap H$ 的普通二次奇点的点 $(x,H)$ 构成的子集为开集.}

注意, 除了 $\operatorname{char}(k)=2$ 且 $\dim X$ 为奇数的情形外, 4.1 均由 3.3 得出.

\subsubsection*{4.1.1.}
设 $R_1$ 为开集 4.1 在 $\mathbb P(N)$ 中的闭补集, 设
\[
F_1=\varphi(R_1)\subset\check X\subset\check{\mathbb P}^{r}
\]
为其像.

\subsubsection*{引理 4.1.2.}
\underline{若开集 4.1 非空, 则 $\operatorname{codim}_{\check{\mathbb P}^{r}}(F_1)\geq2$.}

\noindent\underline{证明.}
由假设,
\[
\dim R_1<\dim\mathbb P(N)=r-1,
\]
并且 $F_1=\varphi(R_1)$.

\subsubsection*{4.1.3.}
设 $F_2$ 为 $\check X$ 中由至少在两个点处与 $X$ 相切的超平面构成的子集. 注意 $F_2$ 不闭, 但有
\begin{equation}\tag{4.1.4}
\overline{F_2}\subset F_1\cup F_2.
\end{equation}
这由 $F_1\cup F_2=F$(参见 3.2)为闭集得出.(包含关系 (4.1.4)“说明”两个奇点的\underline{汇合}不可能为普通二次奇点.) 无论如何, 4.2.4 中的构造将表明 $F_2$ 为可构造集.

为证明 2.5.1, 将应用判别条件 3.2.4. 为此需知
\[
\operatorname{codim}_{\check{\mathbb P}^{r}}(F_1)\geq2
\quad\text{及}\quad
\operatorname{codim}_{\check{\mathbb P}^{r}}(F_2)\geq2.
\]
一方面, 3.7.1 和 4.1.2 表明, 将给定嵌入换为它的第 $d$ 倍嵌入, $d\geq2$, 之后有 $\operatorname{codim}_{\check{\mathbb P}^{r}}(F_1)\geq2$.
\footnote{印本此句写为“3.8.1”和 $F_2$. 普通二次构造为 (3.7.1), 而引理 4.1.2 适用于 $F_1$; 命题 4.2 给出关于 $F_2$ 的所需断言.}
所以, 一旦知道以下命题, 就得到 2.5.1.

\subsubsection*{命题 4.2.}
\underline{设 $X$ 为 $\mathbb P^r$ 中维数 $n\geq1$ 的光滑不可约子簇. 对每个整数 $d\geq2$,}
\[
\check{\mathbb P}^{N}
=\check{\mathbb P}^{\binom{r+d}{d}-1}
\]
\underline{中由至少在两个点处与 $X$ 相切的 $\mathbb P^r$ 中 $d$ 次超曲面所成的子集 $F_2$, 在 $\check{\mathbb P}^{N}$ 中的余维至少为 $2$.}
\footnote{印本此处写为 $\binom{n+d}{d}-1$, 在命题 4.3 中写为 $\binom{r+d}{d}-d$. 2.4 中定义的 $d$ 次 Veronese 嵌入在两处均给出 $N=\binom{r+d}{d}-1$.}